% !PRINTED 251
\weberpage{251}
%<<CONT
If one therefore sets $\gamma - \gamma_0 = \alpha\xi$ (in case $\alpha$ is different from zero), it follows that $\delta - \delta_0 = \beta\xi$, hence
\[ \gamma = \gamma_0 + \alpha\xi, \quad \delta = \delta_0 + \beta\xi, \]
and conversely every pair of numbers $\gamma, \delta$ contained in this form also satisfies it. This evidently holds also in the excepted case $\alpha = 0$. Since we can now put $p^n$ different numbers for $\xi$, the degree of our group is therefore
\weq{4}{p^n (p^{2n} - 1).}

If $p$ is odd, we can derive a still smaller group:

For if $p = 2$, a number $\alpha$ is not different from $-\alpha$, and hence the elements
\weq{5}{\begin{pmatrix} \alpha, & \beta \\ \gamma, & \delta \end{pmatrix}, \quad \begin{pmatrix} -\alpha, & -\beta \\ -\gamma, & -\delta \end{pmatrix}}
are also identical with one another. But if $p$ is odd, the two elements (5) are different from one another, and the whole system of these substitutions splits into pairs of the form (5).

If we take two such pairs and compose each element of the one pair with each element of the other according to the rule (2), we obtain only two different elements, which again form such a pair. These pairs can therefore themselves be regarded as elements of a new group of degree
\weq{6}{\dfrac{p^n (p^{2n} - 1)}{2}}
We can also express this by saying that two elements $A$ whose corresponding numbers differ only in sign are to be regarded as identical.

The group so defined, whose degree is expressed for $p = 2$ by (4) and for an odd $p$ by (6), we shall call \gesperrt{the congruence group belonging to the field $\mathfrak{C}$, and denote it by $E$}.

For the closer investigation of this group it is of importance to know a system of generating substitutions. If, following \S.\,64, (2), we represent the numbers of $\mathfrak{C}$ in the form:
\weq{7}{\xi = x_0 + x_1 \varepsilon_1 + \cdots + x_{n-1} \varepsilon^{n-1},} % [sic] first power printed as subscript \varepsilon_1, not \varepsilon or \varepsilon^1
in which the $x_i$ are rational numbers taken with respect to the modulus $p$
%CONT>>
